% 第五章 服务器系统设计

\section{服务器系统设计}

\subsection{服务器架构设计}
\par 服务器端采用Node.js实现，整体以Express作为Web框架，提供数据接收、鉴权、静态页面托管与控制接口。服务器在启动阶段加载配置文件、历史数据文件与系统设置，并根据配置初始化MQTT管理器以建立与网关的控制链路。该设计使“数据上报链路”和“控制链路”在同一服务内聚合，便于部署与调试。
\par 为降低耦合度并提升可维护性，服务器端按职责划分为数据管理、配置管理、用户管理与MQTT管理等模块。配置管理模块负责读取与合并默认配置与用户配置，统一提供服务器监听地址、数据点上限、自动清理策略与MQTT连接参数等配置项。数据管理模块负责传感器数据的加载、追加、裁剪与统计，并对图表设置与系统设置进行持久化管理。用户管理模块负责用户注册、登录与会话校验，并为需要权限控制的接口提供认证依据。MQTT管理模块负责建立与MQTT Broker的连接、订阅响应主题、发布命令主题，并提供“发送命令并等待响应”的上层能力。

\par 配图建议如下。图5-1建议为“服务器架构模块划分图”，画法为模块框图：展示Express框架、数据管理模块、配置管理模块、用户管理模块、MQTT管理模块，并用连线标注模块间调用关系与数据流向。图5-2建议为“服务器启动与初始化流程示意图”，画法为流程图：加载配置→加载历史数据→加载设置→初始化MQTT→启动HTTP服务，并在关键步骤标注文件路径与初始化结果。
\begin{figure}[H]
\centering
\fbox{\parbox{0.92\linewidth}{\centering 图5-1占位：服务器架构模块划分图。\\
要求展示Express框架与四大管理模块（数据/配置/用户/MQTT），标注模块间调用关系。}}
\caption{服务器架构模块划分图（建议配图）}
\label{fig:server_modules}
\end{figure}
\begin{figure}[H]
\centering
\fbox{\parbox{0.92\linewidth}{\centering 图5-2占位：服务器启动与初始化流程示意图。\\
要求展示从启动到就绪的完整流程，标注配置加载、数据加载、MQTT初始化等关键步骤。}}
\caption{服务器启动与初始化流程示意图（建议配图）}
\label{fig:server_init}
\end{figure}

\subsection{数据管理模块}
\par 数据管理模块以本地JSON文件作为持久化介质，适合毕业设计原型的快速实现与跨平台部署。服务器在接收到网关上传的JSON数据后，会补充服务器时间戳字段并将记录追加到内存数据列表，同时写入数据文件。为避免数据无限增长导致性能下降，模块提供最大数据点数控制并在追加数据后进行裁剪，保留最新的若干条记录以满足可视化需求。
\par 在数据查询方面，服务器提供获取数据列表、获取最新数据、按时间范围查询与导出等接口，便于Web前端绘图与语音助手查询。为提升页面体验，服务器还提供统计接口，能够计算温度、湿度、气压、VOC与气体传感器等指标的最大值、最小值与平均值，用于仪表板摘要展示。图表设置与系统设置通过独立文件管理，使得采样点上限、刷新频率与图表样式等参数可在运行时动态调整并持久化保存。

\par 配图建议如下。图5-3建议为“数据管理模块流程图”，画法为流程图：数据接收→添加时间戳→追加到内存列表→数据裁剪（超过上限）→写入JSON文件，并在裁剪步骤标注保留最新N条的策略。
\begin{figure}[H]
\centering
\fbox{\parbox{0.92\linewidth}{\centering 图5-3占位：数据管理模块流程图。\\
要求展示从数据接收到持久化的完整流程，标注时间戳添加、数据裁剪（保留最新N条）等关键步骤。}}
\caption{数据管理模块流程图（建议配图）}
\label{fig:data_flow_server}
\end{figure}

\subsection{用户认证系统}
\par 用户认证系统用于保护控制接口与敏感数据接口，避免未授权访问。服务器提供注册与登录接口，登录成功后生成会话令牌并通过HTTP响应返回，同时也可写入Cookie以便浏览器端持久化。后续请求在访问需要认证的接口时携带令牌，服务器通过会话校验中间件验证令牌有效性，并将用户信息注入请求上下文。
\par 权限控制采用“默认需要登录”的策略。与数据展示相关的页面（如仪表板、人脸系统页与设置页）在访问时进行登录校验，未登录用户会被重定向到登录页。涉及系统配置修改、数据清空、OpenMV控制命令下发与识别结果查询等接口均需通过认证中间件校验，以降低误操作风险并提升系统安全性基础。

\par 配图建议如下。图5-4建议为“用户认证与会话管理流程示意图”，画法为流程图：用户注册→登录→生成token→写入Cookie/返回响应→后续请求携带token→中间件校验→允许/拒绝访问，并在关键步骤标注token生成与校验逻辑。
\begin{figure}[H]
\centering
\fbox{\parbox{0.92\linewidth}{\centering 图5-4占位：用户认证与会话管理流程示意图。\\
要求展示注册/登录/token生成/中间件校验的完整流程，标注token在Cookie与Authorization头中的使用方式。}}
\caption{用户认证与会话管理流程示意图（建议配图）}
\label{fig:auth_flow}
\end{figure}

\subsection{Web前端设计}
\par Web前端采用静态页面方式部署，由服务器端统一托管。页面结构围绕三类核心场景组织，分别为环境数据仪表板、人脸系统管理与系统设置。仪表板页面在加载后首先校验登录状态，然后并行拉取传感器数据与统计信息，使用Chart.js绘制多参数曲线，并在页面顶部展示最后更新时间与连接状态。设置页面用于修改最大数据点数、刷新频率与图表样式等参数，并通过调用服务器API完成保存与生效。人脸系统页面提供人脸采集、识别、停止、图像获取、名单查询与删除等操作入口，并展示实时图像与识别结果列表。
\par 数据更新策略以轮询为主。仪表板按配置的刷新周期拉取最新数据并增量更新图表，从而在无需额外长连接的情况下实现“准实时”展示。人脸系统页面除命令触发外，还会周期拉取识别结果列表，并根据图像更新频率请求最新图像并渲染，同时在浏览器侧叠加检测框以增强可解释性。考虑到毕业设计部署环境的简化，本系统未强依赖WebSocket推送，但整体架构保留了后续升级为实时推送的空间。

\par 配图建议如下。图5-5建议为“Web前端页面结构图”，画法为页面布局图：展示仪表板页面（图表区域、统计摘要、状态栏）、人脸系统页面（图像区域、控制按钮、识别结果列表）、设置页面（参数配置表单、图表样式设置），并用标注说明各区域功能。
\begin{figure}[H]
\centering
\fbox{\parbox{0.92\linewidth}{\centering 图5-5占位：Web前端页面结构图。\\
要求展示仪表板/人脸系统/设置三个核心页面的布局结构，标注各功能区域（图表、按钮、列表、表单等）。}}
\caption{Web前端页面结构图（建议配图）}
\label{fig:web_structure}
\end{figure}

\subsection{设备接入与管理模块}
\par 设备接入与管理在服务器侧主要体现为两部分。一部分是HTTP数据接收接口，用于接收网关周期上报的环境数据并写入存储。另一部分是MQTT控制管理，用于向网关下发OpenMV相关控制命令并等待响应。服务器在初始化时依据配置文件中的MQTT参数建立连接并订阅响应主题，收到响应后根据响应内容类型进行解析与分发。对于识别结果等异步消息，服务器会将其保存为可查询的历史记录，便于前端持续展示。
\par 在控制接口设计上，服务器对OpenMV命令集进行白名单校验，避免任意字符串被直接下发造成不可控行为。服务器提供“发送命令并等待响应”的封装能力，允许前端指定超时时间以适配不同命令的耗时特性。对于图像数据，服务器与前端共同约定以Base64字符串作为载体，并在必要时支持分块消息的拼装，从而在MQTT缓冲区限制下仍可完成图像回传与显示。
\par 在线状态与可维护性方面，服务器可通过MQTT连接状态与接口调用成功率对设备可用性进行间接判断。作为毕业设计原型，本文以实现链路闭环与可用性验证为主，未引入专门的心跳协议与设备身份体系；在后续扩展中，可进一步加入设备ID与密钥管理、心跳与离线检测、以及配置下发与固件升级通知等能力，以增强系统的工程完整性。

\par 配图建议如下。图5-6建议为“设备接入与控制接口关系图”，画法为接口关系图：展示HTTP数据接收接口、MQTT命令下发接口、前端API接口（/data、/api/openmv/command等），并用连线标注接口调用关系与数据流向。
\begin{figure}[H]
\centering
\fbox{\parbox{0.92\linewidth}{\centering 图5-6占位：设备接入与控制接口关系图。\\
要求展示HTTP数据接收、MQTT控制、前端API等接口，标注接口路径、请求方法（GET/POST）与数据流向。}}
\caption{设备接入与控制接口关系图（建议配图）}
\label{fig:api_structure}
\end{figure}
